\section{Conclusion}\label{sec:conclusion}
This study asked how news sentiment can improve the optimised allocation of a weekly, long-only U.S. and euro-area equity portfolio after costs, with all models chosen on 2012--2014 and tested once on 2015--2025.

\textbf{RQ1.} Sentiment signals mostly followed price moves; the small continuation over the next month faded by day 40 (H1a in direction only). Firm sentiment had a small, insignificant positive effect after price controls (H1b in direction only).

\textbf{RQ2.} Correcting the factor forecast with sentiment did not raise its IC, but the portfolio earned \BootCombinedPriceCAGR{} percentage points more a year with a smaller drawdown, a gain earned on both buy and sell decisions (H2a--H2c in direction only).

\textbf{RQ3.} The sentiment-driven allocation earned \SentimentCAGR{} a year, slightly below equal weighting, with a Sharpe ratio of \SentimentSharpe{}, marginally above equal weighting (\EWSharpe), minimum variance (\MinVarSharpe) and MSCI World (\IndexSharpe), and a smaller drawdown than equal weighting and the index. The differences are not significant and depend on costs and period (H3 in direction only).

Sentiment improved the allocation in direction without improving the return forecasts, and no improvement was statistically significant. Gains of one to two percentage points a year are small, but in liquid large caps they are economically relevant. For a portfolio manager, public headline sentiment is useful as a confirmation layer for a price-based allocation and weak as a source of alpha on its own.
